\section{Phê bình và giới hạn}
\begin{frame}{Điều paper làm tốt}
  \begin{itemize}
    \item Đặt đúng giao điểm của ba vấn đề: đa phương thức, tuần tự và MLLM.
    \item Chuyển image thành text để giải quyết context dài bằng một cơ chế đơn giản, dễ triển khai.
    \item Recurrent prompt tạo ra intermediate preference có thể đọc và phân tích.
    \item So sánh đa dataset, nhiều nhóm baseline, có ablation và parameter analysis.
  \end{itemize}
\end{frame}

\begin{frame}{Các giới hạn cần hỏi thêm}
  \begin{itemize}
    \item \textbf{Information bottleneck:} mọi visual cue cuối cùng bị nén qua generated text summary.
    \item \textbf{Chi phí inference:} mỗi block cần gọi LLM; history dài làm tăng số lần gọi và độ trễ.
    \item \textbf{Evaluation gap:} paper chứng minh offline metrics, chưa chứng minh latency, cost hoặc online user satisfaction.
    \item \textbf{Reproducibility gap:} code checkout hiện tại có hyperparameter khác paper, nên cần log rõ cấu hình để tái hiện.
    \item \textbf{Calibration:} xác suất yes/(yes+no) phụ thuộc tokenization và calibration của MLLM, chưa phải xác suất tương tác tuyệt đối.
  \end{itemize}
\end{frame}

